\documentclass[11pt]{article}
\usepackage[utf8]{inputenc}
\usepackage[T2A]{fontenc}
\usepackage[english,russian]{babel}
\usepackage{amsmath,amssymb,amsthm}
\usepackage{booktabs}
\usepackage{geometry}
\usepackage{hyperref}
\usepackage{longtable}
\usepackage{natbib}
\usepackage{url}
\geometry{margin=1in}

\title{Квинтет IDEAV: компилируемая метаданными модель данных для low-code бизнес-приложений}
\author{Алексей П. Семенов\\Integram / IDEAV}
\date{Черновик для arXiv cs.DB, май 2026}

\newtheorem{definition}{Определение}
\newtheorem{theorem}{Теорема}

\begin{document}
\maketitle

\begin{abstract}
Схемы Entity--Attribute--Value (EAV) обычно рассматриваются как необходимый, но
опасный компромисс: они позволяют расширять прикладную схему без миграций, но
переносят типизацию, целостность, навигацию и планирование запросов из
реляционной модели в неформальный код приложения. В статье описывается
квинтет IDEAV, компилируемая метаданными модель данных, используемая в
low-code платформе Integram. Цель модели состоит в том, чтобы сохранить
пластичность EAV, но сделать недостающую семантику явной. Все прикладные
метаданные и значения представлены в едином пространстве кортежей, однако их
роли разделены на пять взаимодействующих элементов: типы, реквизиты, объекты,
связи и компилятор. Компилятор, реализованный процедурой
\texttt{Compile\_Report}, преобразует пользовательское описание отчета в SQL,
следуя четырем формально заданным случаям навигации: ссылка текущего объекта на
новый тип, дочерний массив текущего объекта, ссылка нового типа обратно на
текущий объект и переход от дочернего объекта к его родительскому массиву. Мы
даем формальное описание модели, набросок доказательства того, что компилятор
покрывает любой связный граф отчета, выразимый метаданными, и обсуждаем
физическое хранение с опорой на патентное семейство RU2650032C1 и
US11138174B2. Сравнение с медицинскими EAV-системами, коммерческими
антипаттернами и современными low-code платформами показывает, что проблема EAV
заключается не в строковом хранении атрибутов как таковом, а в отсутствии
формальной модели и компилятора. Так как открытые контролируемые бенчмарки пока
не опубликованы, утверждения о производительности сформулированы как гипотезы
и план измерений, а не как готовые экспериментальные результаты.
\end{abstract}

\noindent\textbf{Ключевые слова:} low-code, метаданные, EAV, эволюция схемы, компиляция запросов.

\section{Введение}

Low-code платформа обещает, что бизнес-аналитик сможет изменить процесс без
ожидания миграции базы данных. Классическое реляционное проектирование обещает
другую ценность: структуры данных заранее объявлены, типизированы, нормализованы
и оптимизируемы \citep{codd1970relational}. Напряжение между этими обещаниями
и объясняет постоянное возвращение схем Entity--Attribute--Value. EAV легко
расширять, но неструктурированный EAV часто превращается в динамическое
строковое хранилище, в котором СУБД уже не видит реальную схему
\citep{karwin2010sql,kyte2003effective}.

Поэтому главный вопрос звучит так: почему не взять обычный EAV? Потому что EAV
сам по себе не является моделью. Это физический прием кодирования. Он говорит,
что атрибуты могут быть строками, но не говорит, как атрибуты типизируются, как
отличать ссылку от скалярного значения, как восстанавливать соединения, как
представлять массивы и какая инвариантная логика не дает набору троек стать
частным языком конкретного приложения.

IDEAV исходит из той же потребности в позднем изменении структуры, но не
останавливается на строках атрибутов. Центральная идея модели --- квинтет:

\[
  \mathcal{Q}_{IDEAV} =
  \langle \mathcal{T}, \mathcal{R}, \mathcal{O}, \mathcal{L}, \mathcal{C} \rangle ,
\]

где \(\mathcal{T}\) --- множество логических типов, \(\mathcal{R}\) ---
множество реквизитов, \(\mathcal{O}\) --- множество объектов, \(\mathcal{L}\)
--- множество типизированных связей, а \(\mathcal{C}\) --- компилятор
метаданных. Пятый элемент принципиален. Без компилятора первые четыре элемента
рискуют остаться переименованным EAV. С компилятором отчеты, формы и API
работают не с произвольными ключами и строками, а с графом объявленных
метаданных.

Модель развивалась в коммерческой платформе Integram с 2015 года; концептуальная
работа началась в 2003 году. Она также связана с патентным семейством
``Electronic database and method of its formation'' \citep{semenov2018ru} и
``Electronic database and method for forming same'' \citep{semenov2021us}.
Патенты здесь не используются как рекламный аргумент. Они фиксируют
документированную линию разработки однотабличного представления данных и
метаданных с кодированием ролей.

\section{Связанные работы}

\subsection{Медицинский и научный EAV}

Самая сильная академическая традиция EAV возникла не в коммерции, а в медицине
и научном сборе данных. ACT/DB применял attribute-value хранение для переменных
клинических исследований, структура которых менялась во времени
\citep{nadkarni1998actdb}. EAV/CR дополнил идею классами и отношениями для
гетерогенных научных данных \citep{nadkarni1999eavcr,marenco2003eavcr}.
Веб-базы клинических данных и i2b2 показали, что онтологические метаданные могут
поддерживать богатые развивающиеся наборы данных
\citep{anhoj2003generic,murphy2010i2b2}.

Эти системы частично отвечают на вопрос ``почему не EAV'': EAV работает, когда
предметная область признает метаданные первоклассным объектом. Но они также
показывают недостающее звено для операционного бизнес-софта. Исследовательская
медицинская схема часто оптимизирована под разреженные наблюдения и последующий
анализ. CRM или ERP дополнительно требуют ежедневных форм, прав доступа, ссылок,
дочерних таблиц, импорта и пользовательских отчетов. Поэтому IDEAV трактует
метаданные не только как описание, но и как исполняемую схему для компилятора.

\subsection{Коммерческие применения и типовые сбои}

Бизнес-системы применяли EAV для характеристик товаров, пользовательских полей и
CRM-расширений. Предупреждающая литература хорошо известна: статья ``Bad
CaRMa'' описывает боль каталожной системы, в которой entity-attribute дизайн
скрыл полезную структуру от СУБД \citep{gorman2006badcarma}. Обсуждения SQL
антипаттернов формулируют похожие выводы: универсальные строки атрибутов
ослабляют ограничения, усложняют запросы и выталкивают целостность в приложение
\citep{karwin2010sql}.

Magento и WooCommerce дают более тонкий пример. Adobe Commerce продолжает
документировать расширяемые атрибуты \citep{adobeCommerceAttributes}, а
WooCommerce ввел таблицы поиска товаров и High-Performance Order Storage, чтобы
снизить зависимость горячих путей от широких meta-таблиц
\citep{woocommerceLookup,woocommerceHPOS}. Это не просто ``провал EAV''. Это
свидетельство того, что кодирование расширяемости без формального компилятора,
типизированных правил навигации и проекций под нагрузку позже требует вторичных
структур.

В черновике сознательно не используются как первичное доказательство
неархивированные истории, включая пересказы о cascade-delete инцидентах в CRM
продуктах. Для академического текста сильнее опираться на публичные источники и
проверяемые свойства модели.
По той же причине сравнение с Kustomer оставлено только как кандидат для будущей
архивной проверки: без первичного публичного источника оно не должно нести
эмпирическое утверждение.

\subsection{Современные low-code платформы}

Low-code системы решают эволюцию схемы разными способами. Одни дают
таблично-спредшитный интерфейс, другие подключаются к внешней реляционной БД,
третьи генерируют обычные таблицы, четвертые держат метаданные внутри рантайма.
Таблица~\ref{tab:lowcode-ru} сравнивает их на уровне модели.

\begin{longtable}{p{0.18\linewidth}p{0.31\linewidth}p{0.39\linewidth}}
\caption{Выбранные low-code системы и стратегия схемы.}\label{tab:lowcode-ru}\\
\toprule
Система & Стратегия хранения & Следствие для вопроса об EAV \\
\midrule
\endfirsthead
\toprule
Система & Стратегия хранения & Следствие для вопроса об EAV \\
\midrule
\endhead
Airtable & Hosted base/table abstraction with API access \citep{airtableApi} &
Сильный пользовательский опыт, но семантика хранения и компиляции проприетарна. \\
Bubble & Прикладные data types в рантайме Bubble \citep{bubbleData} &
Гибко, но модель БД не является переносимым формальным контрактом. \\
Retool & Подключает формы и workflow к внешним ресурсам
\citep{retoolResources} & Оставляет настоящие БД доступными; эволюция схемы
делегируется подключенным системам. \\
NocoDB, Baserow & Spreadsheet-like интерфейсы над структурированными backend
\citep{nocodb,baserow} & Удобный low-code слой, обычно таблично-ориентированный,
а не единое формальное пространство кортежей. \\
Directus, Strapi & Headless data/content платформы поверх схем БД
\citep{directus,strapi} & Хороши, когда схема уже задана; менее сфокусированы на
компиляции произвольных бизнес-метаданных в соединения. \\
Caspio, Knack, AppSheet & Hosted app builders \citep{caspio,knack,appsheet} &
Снижают стоимость разработки, но не публикуют формальный аналог компилятора
IDEAV. \\
Integram / IDEAV & Единое пространство кортежей с кодированием ролей плюс
компилятор метаданных \citep{integramWorkflow2026,integramMcp2026} &
Рассматривает эволюцию схемы и навигацию запросов как одну формальную задачу. \\
\bottomrule
\end{longtable}

\subsection{Контекст систем баз данных}

IDEAV не является заменой реляционной алгебры, NewSQL или document stores. Это
метамодель более высокого уровня, использующая реляционную СУБД как substrate
исполнения. NewSQL показывает, что реляционные системы масштабируют
транзакционные нагрузки, когда схема известна \citep{pavlo2016newsql}. NoSQL и
polyglot persistence показывают, что альтернативные модели упрощают отдельные
паттерны доступа \citep{sadalage2012nosql,floratou2012nosql}. IDEAV задает
более узкий вопрос: если бизнес-пользователи создают схему во время работы, какие
формальные метаданные нужны, чтобы сгенерированный реляционный запрос оставался
предсказуемым?

\section{Формальная модель квинтета IDEAV}

\subsection{Физическое пространство кортежей}

Патентное ядро описывает электронную базу данных, элементы которой организованы
в одной таблице с идентификаторами, родительскими идентификаторами, значениями и
кодом типа \citep{semenov2018ru,semenov2021us}. Реализация Integram расширяет
этот минимальный кортеж служебными колонками метаданных, используемыми
компилятором. Физическое множество можно записать так:

\[
  U = \{q_i = (id_i, t_i, up_i, val_i, A_i)\},
\]

где \(id_i\) --- стабильный идентификатор, \(t_i\) --- идентификатор типа,
\(up_i\) --- родительский или объектный идентификатор, \(val_i\) --- значение,
а \(A_i\) --- конечная карта служебных атрибутов: aliases, маркеры ссылок,
маркеры массивов, признаки уникальности и права доступа. Это не ``просто EAV'',
потому что \(t\) и \(up\) интерпретируются объявленными ролями метаданных, а не
каждым экраном приложения по-своему.

\subsection{Пять ролей}

\begin{definition}[Тип]
Тип \(\tau \in \mathcal{T}\) объявляет логический бизнес-класс. В API он
создается командой \texttt{\_d\_new}. Первое поле имени может адресоваться как
\texttt{t\(\tau\)}.
\[
  \tau = (id, name, parent, unique, permissions).
\]
\end{definition}

\begin{definition}[Реквизит]
Реквизит \(r \in \mathcal{R}\) --- поле или структурное отношение, принадлежащее
типу. Он создается командой \texttt{\_d\_req}. Скалярный реквизит хранит
значения; ссылочный содержит \texttt{ref} или \texttt{ref\_id}; массивный
содержит \texttt{arr\_id}.
\[
  r = (id, owner, name, domain, ref, ref\_id, arr\_id, attrs).
\]
\end{definition}

\begin{definition}[Объект]
Объект \(o \in \mathcal{O}\) --- типизированная бизнес-запись.
\[
  o = (id, type(o), parent(o), name(o), \{r \mapsto v_r\}).
\]
Скалярные значения хранятся кортежами под объектом; имя записи хранится как
значение строки самого типа.
\end{definition}

\begin{definition}[Связь]
Связь \(\ell \in \mathcal{L}\) --- типизированный факт навигации между
объектами или между родителем и дочерним элементом массива. Ссылки и массивы
разделены, потому что они компилируются в разные предикаты соединения.
\[
  \ell_{ref}: o_s \xrightarrow{r} o_t,\qquad
  \ell_{arr}: o_p \xrightarrow{a} o_c.
\]
\end{definition}

\begin{definition}[Компилятор]
Компилятор \(\mathcal{C}\) отображает метаданные и определение отчета в SQL:
\[
  \mathcal{C}: (\mathcal{T},\mathcal{R},\mathcal{L}, Report, \Gamma)
  \rightarrow SQL ,
\]
где \(\Gamma\) содержит aliases, фильтры, функции и текущий контекст блока. В
реализации эту роль выполняет \texttt{Compile\_Report}.
\end{definition}

\subsection{Компилятор метаданных}

Отчет --- это список требуемых колонок. Каждая колонка принадлежит логическому
типу. Компилятор поддерживает множество \(J\) логических типов, уже подключенных
к SQL-запросу. Для каждого нового требуемого типа \(n\) он ищет тип
\(j \in J\), который может быть соединен с \(n\). Документация и код выделяют
четыре случая:

\begin{center}
\begin{tabular}{lll}
\toprule
Случай & Условие метаданных & Форма соединения \\
\midrule
Reference to us & \(j\) имеет ссылочный реквизит на \(n\) &
\(j.ref = n.id\) \\
We are an Array & \(n\) является дочерним массивом \(j\) &
\(n.up = j.id\) \\
We have a Reference & \(n\) имеет ссылочный реквизит на \(j\) &
\(n.ref = j.id\) \\
We got an Array & \(j\) является дочерним массивом \(n\) &
\(j.up = n.id\) \\
\bottomrule
\end{tabular}
\end{center}

Упрощенный алгоритм:

\begin{verbatim}
J := {master type from the first report column}
while uncompiled columns remain:
    progress := false
    for column c with type n not in J:
        for each j in J:
            if j references n: add join "j.ref = n.id"
            else if n is array child of j: add join "n.up = j.id"
            else if n references j: add join "n.ref = j.id"
            else if j is array child of n: add join "j.up = n.id"
            if a case matched:
                J := J union {n}
                progress := true
                break
    if not progress: fail with "cannot connect report columns"
\end{verbatim}

Здесь IDEAV расходится с обычным EAV. Универсальная таблица атрибутов может
хранить ``customer'', ``invoice'' и ``invoice lines'' как строки, но она не
обязана знать, какое соединение допустимо. IDEAV делает четыре навигации явными
метаданными и требует, чтобы компилятор выбрал одну из них.

\subsection{Набросок корректности}

\begin{theorem}[Компиляция связного отчета]
Пусть \(G_M=(V,E)\) --- неориентированный граф, индуцированный IDEAV
метаданными ссылок и массивов для отчета: вершины являются логическими типами, а
ребра --- допустимыми экземплярами одного из четырех случаев компилятора. Если
требуемые типы отчета образуют связный подграф \(G_M\), и каждое ребро имеет
однозначную роль метаданных, то \(\mathcal{C}\) может добавить путь соединений
для каждого требуемого типа.
\end{theorem}

\begin{proof}
Выберем первый тип отчета как master vertex. Инвариант после каждого успешного
шага компилятора: \(J\) является связным множеством вершин, а сгенерированный SQL
содержит путь от master до каждой вершины из \(J\). В начале инвариант верен
для singleton master. Пусть он верен для \(k\) вершин. Так как требуемый подграф
связен, если не все вершины входят в \(J\), существует ребро из некоторого
\(j \in J\) в некоторый \(n \notin J\). По построению каждое ребро метаданных
является ровно одним из четырех случаев компилятора, поэтому компилятор может
добавить соответствующий join predicate и расширить \(J\) до
\(J \cup \{n\}\). Инвариант сохраняется. По индукции по числу требуемых типов
все вершины будут соединены. Если такого ребра нет, отчет несвязен или
неоднозначен, и отказ является корректным поведением.
\end{proof}

Теорема не утверждает оптимальность стоимости. Она утверждает семантическое
покрытие для графа метаданных, который сама платформа позволяет пользователям
создавать. Стоимость запроса относится к движку исполнения и нагрузке, а не к
следствию доказательства.

\subsection{Физическое хранение и индексы}

Дизайн хранения держит метаданные и данные в одной физической relation, обычно
с составными индексами по identifiers, type, parent и value
\citep{integramWorkflow2026}. Фиксированная форма кортежа дает два инженерных
преимущества: изменение схемы становится изменением данных, а компилятор
нацеливается на стабильный набор индексируемых колонок. Но было бы неверно
утверждать, что это само по себе гарантирует постоянную производительность.
Высота B-tree зависит от числа строк, размера страницы, распределения ключей и
стратегии партиционирования. Защищаемое утверждение более узкое: фиксированная
tuple schema и partitionable identifier ranges дают реализации устойчивый
механизм удержания операций с метаданными и сгенерированных соединений в
известном эксплуатационном диапазоне. Контролируемые бенчмарки нужны, чтобы
превратить этот механизм в измеренные latency claims.

\section{Эмпирическая валидация}

Модель нужно проверять на трех уровнях.

\paragraph{Корректность компилятора.}
Синтетические графы метаданных должны покрывать каждую допустимую комбинацию
ссылок и массивов. Для каждого графа сгенерированный отчет следует сравнить с
эквивалентным вручную написанным relational query над нормализованной схемой.
Минимальный набор тестов содержит chains, stars, diamonds, reverse references,
parent-to-child arrays, child-to-parent arrays и несвязные отчеты, которые
обязаны завершаться ошибкой.

\paragraph{Эволюция схемы.}
Нагрузка должна создавать новые типы, добавлять скалярные реквизиты, ссылки и
массивы, а затем выполнять отчеты без DDL migrations. Сравнивать следует с
обычными migrations, generic EAV table, JSON/document column design и
metadata-compiled IDEAV design.

\paragraph{Производительность.}
Бенчмарки должны публиковать latency записи, время компиляции отчета, время
исполнения сгенерированного SQL, размер индексов и storage overhead как функции:

\[
  N_o = \text{objects},\quad
  N_t = \text{types},\quad
  N_r = \text{requisites per type},\quad
  d = \text{report join depth}.
\]

Пока такие результаты не опубликованы, промышленная история с 2015 года должна
рассматриваться как свидетельство реализуемости, но не как контролируемое
доказательство производительности. Поэтому сам протокол измерений является
частью вклада: он показывает, что именно надо измерить перед честным сравнением
IDEAV с relational, EAV и JSON alternatives.

\section{Обсуждение}

IDEAV меняет постановку спора об EAV. Проблема EAV в бизнес-системах не в том,
что атрибуты являются строками. Проблема в том, что строковое кодирование часто
заставляют одновременно нести схему, ограничения, навигацию, поведение UI и
семантику отчетов без формального слоя. IDEAV назначает эти обязанности
метаданным и делает метаданные исполняемыми через компилятор.

Это также меняет роль AI-assisted development. Языковая модель или scripted
agent может вызывать компактный API: \texttt{\_d\_new}, \texttt{\_d\_req},
\texttt{\_d\_ref}, \texttt{\_m\_new}, \texttt{\_m\_save}. Эти вызовы не просто
создают таблицы в UI; они расширяют тот же граф метаданных, который позднее
используется компилятором отчетов и форм. AI surface оказывается меньше, чем
полная поверхность SQL migration, но формальнее, чем screen-scraping automation.

Ограничения существенны. Сгенерированные соединения могут быть семантически
корректными и при этом медленными. Аналитические нагрузки могут требовать
materialized projections. Строгая типизация, история миграций и multi-tenant
изоляция должны быть специфицированы отдельно. Модель также зависит от
качества управления метаданными: если пользователи создают плохие типы и связи,
компилятор сохранит их смысл, но не сделает предметную область хорошо
спроектированной.

\section{Заключение}

EAV часто терпел неудачу в бизнес-системах потому, что использовался как обход
моделирования, а не как слой внутри формальной модели. Квинтет IDEAV предлагает
другой контракт: сохранить расширяемое пространство кортежей, но потребовать
явные типы, реквизиты, объекты, связи и компилятор. В результате получается
low-code модель данных, где эволюция схемы дешева, но не невидима; отношения
динамичны, но не произвольны; а сгенерированный SQL следует правилам
метаданных, которые можно формулировать, тестировать и ревьюировать.

Следующий исследовательский шаг --- эмпирический. Открытые бенчмарки должны
сравнить IDEAV с классическим EAV, нормализованными реляционными схемами и
JSON/document designs по эволюции схемы, компиляции отчетов и исполнению
запросов. Если измерения подтвердят операционные гипотезы, IDEAV следует
рассматривать не как отрицание реляционных баз данных, а как дисциплинированный
метаданный слой над ними.

\appendix

\section{Патентные ссылки}

RU2650032C1, ``Electronic database and method of its formation'', указывает
Алексея Петровича Семенова как изобретателя и описывает базу данных, элементы
которой организованы в одной таблице с identifiers, parent identifiers, values
и datatype codes \citep{semenov2018ru}. Дата приоритета --- 2017-03-20, дата
публикации --- 2018-04-06.

US11138174B2, ``Electronic database and method for forming same'', принадлежит
тому же патентному семейству и описывает однотабличную организацию со строками
для root element, datatype, term, term attribute и data \citep{semenov2021us}.
Дата приоритета --- 2017-03-20, дата публикации --- 2021-10-05.

Патенты фиксируют claims новизны вокруг физической организации. Отдельное
утверждение этой статьи находится на уровне модели: IDEAV нужно оценивать как
комбинацию физического tuple space, формальных ролей метаданных и компилятора.

\section{Эскиз API IDEAV}

Документация репозитория описывает компактный API для создания метаданных и
данных \citep{integramMcp2026}. Примеры ниже схематичны:

\begin{verbatim}
# Create a logical type/table.
POST _d_new
  val = "Customer"

# Add a scalar requisite/column to a type.
POST _d_req/{customerTypeId}
  val = "Phone"
  t   = "{stringTypeId}"

# Create a reference domain and attach it as a requisite.
POST _d_ref/{targetTypeId}
POST _d_req/{sourceTypeId}
  val = "Customer"
  t   = "{createdReferenceTypeId}"

# Attach an existing type as a subordinate array.
POST _d_req/{parentTypeId}
  val = "Invoice lines"
  t   = "{lineTypeId}"

# Create and save objects.
POST _m_new/{typeId}
POST _m_save/{objectId}
  t{typeId} = "Object name"
  {requisiteId} = "value"

# Inspect metadata.
GET metadata?JSON=1
\end{verbatim}

Эти команды важны потому, что редактирование схемы в IDEAV не является
отдельной migration subsystem. API изменяет тот же граф метаданных, который
читает компилятор отчетов.

\bibliographystyle{plainnat}
\bibliography{references}

\end{document}
